\documentclass[11pt]{article}
\usepackage[margin=2.5cm]{geometry}
\usepackage{amsmath,amssymb,amsthm}
\usepackage{booktabs}
\usepackage{array}
\usepackage{hyperref}

\newtheorem{example}{Example}
\newcommand{\PY}{\mathrm{PY}}
\newcommand{\Unif}{\mathrm{Uniform}}

\title{Priors over model forms: a description of each model}
\author{}
\date{}

\begin{document}
\maketitle

\noindent This note describes every prior over model forms that we compare. It assumes probability and
statistics at the level of a late-stage undergraduate: conditional probability, the multinomial and
Dirichlet distributions, and the idea of a Markov chain Monte Carlo (MCMC) sampler. Everything else is
built up from there. Section~\ref{sec:setup} sets up the problem; Section~\ref{sec:trees} turns a model
form into a tree and a tree probability into a product of simple choices;
Sections~\ref{sec:smoothing}--\ref{sec:hpy} build the hierarchical Pitman--Yor model that underlies most
of the priors; Sections~\ref{sec:variants}--\ref{sec:pcfg} describe each prior; and
Section~\ref{sec:evaluation} says how they are compared.

%=========================================================================
\section{The problem}\label{sec:setup}

A \emph{model form} is a mathematical expression in one variable $x$ whose constants are left free, for
example the Michaelis--Menten form
\[
  f(x) = \frac{c_0\, x}{c_1 + x}.
\]
Two expressions are the same form if they differ only in the values of their constants. We write every
form in a canonical way (a fixed simplification and ordering of terms, with constants renamed
$c_0, c_1, \dots$ in order of appearance), so that each form has one string, its \emph{key}. Numerical
literals are also replaced by free constants (so $2x$ and $0.693x$ are both the form $c_0 x$), except
integer and half-integer exponents, which are part of the structure (so $x^2$ and $\sqrt{x}$ stay as
they are).

A \emph{prior over model forms} is a probability distribution $P$ over the (infinite) set of all forms.
We want $P$ to put high probability on the forms that biologists and ecologists actually write down, so
we \emph{learn} it from a training corpus: a list of forms extracted from published models in several
sources (BioModels, CellML, ModelDB, SABIO-RK, and so on). Each published model contributes each of
its distinct forms once.

Two practical requirements shape everything below.
\begin{enumerate}
  \item The prior must give positive probability to forms that are \emph{not} in the corpus, because
  the true form of a new system may be one nobody has published.
  \item The prior must be \emph{proper}: the probabilities of all forms must add up to one. Otherwise
  comparisons between priors are meaningless (a prior that gives every form probability $1$ would ``win''
  every comparison).
\end{enumerate}
The simplest idea, the empirical distribution of corpus forms, fails requirement~1: it gives probability
zero to every unseen form. All of the priors below instead build a form out of small pieces, and give
probabilities to the pieces.

%=========================================================================
\section{Forms as trees, and tree probabilities as products}\label{sec:trees}

\subsection{The tree coding}

Every form is written as an \emph{expression tree}. The internal nodes are operators and the leaves are
the variable or a constant. We use the coding of Bartlett, Desmond and Ferreira's \texttt{fnprior}
(the symbol set they call the basis), with 13 symbols:
\begin{center}
\begin{tabular}{lll}
  \toprule
  kind & symbols & number of children (arity)\\
  \midrule
  leaves & \texttt{a} (any constant), \texttt{x} (the variable) & 0\\
  unary operators & \texttt{sqrt}, \texttt{exp}, \texttt{log}, \texttt{sin}, \texttt{tanh}, \texttt{inv} ($1/\cdot$) & 1\\
  binary operators & \texttt{+}, \texttt{-}, \texttt{*}, \texttt{/}, \texttt{pow} & 2\\
  \bottomrule
\end{tabular}
\end{center}
Every constant is coded by the same symbol \texttt{a}: the prior is over forms, not over constant
values.

\begin{example}\label{ex:mm}
The Michaelis--Menten form $c_0 x/(c_1+x)$ is the tree
\[
  \texttt{*}\big(\texttt{a},\ \texttt{/}(\texttt{x},\ \texttt{+}(\texttt{a},\texttt{x}))\big),
\]
which, read from the top down and left to right (\emph{prefix order}), is the sequence of seven symbols
\texttt{* a / x + a x}. Because each symbol's arity is known, this sequence determines the tree
uniquely.
\end{example}

\subsection{A tree probability as a product of choices}

Imagine generating a tree from the top down. We choose the root's symbol. If it has children, we then
choose the first child's symbol, then (recursively) everything below the first child, then the second
child's symbol, and so on. The tree is finished when every branch has ended in a leaf. Write the nodes
in this order as $1, 2, \dots, m$, with symbol $s_i$ at node $i$. By the chain rule of probability,
\begin{equation}\label{eq:chain}
  P(T) = \prod_{i=1}^{m} P(s_i \mid s_1, \dots, s_{i-1}).
\end{equation}
This is exact, but each factor depends on the whole history, which we could never estimate. Every tree
prior here makes the same kind of simplifying assumption: \emph{the choice at node $i$ depends on the
history only through a short summary $u_i$, its context.} Then
\begin{equation}\label{eq:product}
  P(T) = \prod_{i=1}^{m} G_{u_i}(s_i),
\end{equation}
where, for each context $u$, $G_u$ is a probability distribution over the 13 symbols. A prior is then
specified by (i) what goes into the context, and (ii) how the distributions $G_u$ are estimated from the
corpus.

\subsection{Contexts}

The context of a node is built from a few named components:
\begin{itemize}
  \item $a_1$ (the parent's symbol), $a_2$ (the grandparent's), $a_3$ (the great-grandparent's); at the
  root these are all a special ``none'' symbol, so the root has its own context;
  \item $l$, the left sibling's symbol, for a second child;
  \item $d$, the node's depth in the tree (the root has depth $0$), capped at 6;
  \item $s$, the corpus source of the training form (used by one variant only).
\end{itemize}
First children (and the root) and second children are always kept apart: they have separate
distributions even when their components agree. Following \texttt{fnprior}, with ``order'' $n$ a
first child's context is its $n$ nearest ancestors $(a_n,\dots,a_1)$ and a second child's context is its
$n-1$ nearest ancestors and its left sibling, $(a_{n-1},\dots,a_1,l)$. The left sibling matters: in
$x - x$ or $x \cdot x$ the second argument is more predictable once the first is known.

\begin{example}\label{ex:contexts}
With $n=2$ the seven nodes of the Michaelis--Menten tree have these contexts:
\begin{center}
\begin{tabular}{clll}
  \toprule
  node & symbol & kind & context\\
  \midrule
  1 & \texttt{*} & root & (none, none)\\
  2 & \texttt{a} & first child & (grandparent none, parent \texttt{*})\\
  3 & \texttt{/} & second child & (parent \texttt{*}, left sibling \texttt{a})\\
  4 & \texttt{x} & first child & (grandparent \texttt{*}, parent \texttt{/})\\
  5 & \texttt{+} & second child & (parent \texttt{/}, left sibling \texttt{x})\\
  6 & \texttt{a} & first child & (grandparent \texttt{/}, parent \texttt{+})\\
  7 & \texttt{x} & second child & (parent \texttt{+}, left sibling \texttt{a})\\
  \bottomrule
\end{tabular}
\end{center}
so $P(T) = G_{\text{root}}(\texttt{*})\,G_{(\cdot,\texttt{*})}(\texttt{a})\,G_{(\texttt{*},\texttt{a})}(\texttt{/})
\cdots G_{(\texttt{+},\texttt{a})}(\texttt{x})$.
\end{example}

\subsection{When is a tree prior proper?}\label{sec:proper}

Equation~\eqref{eq:product} defines a probability for every finite tree, but the total over all finite
trees need not be one: the generating process can, with positive probability, go on forever. A useful
way to see this is through the theory of branching processes. Each node ``gives birth'' to $0$, $1$ or
$2$ children according to its symbol. If the expected number of children per node is above one, the
tree grows without end with positive probability, and that probability is missing from the finite
trees. For example, choosing every symbol uniformly at random (probability $1/13$ each) gives an
expected $(2\cdot 0 + 6\cdot 1 + 5\cdot 2)/13 = 16/13 > 1$ children per node; we found that only about
60\% of this process's mass is on finite trees, so it is \emph{not} a proper prior. A prior learned from
a corpus of finite trees puts enough mass on leaves that its expected number of children is below one,
and we check numerically, for every prior and every training set, that its mass on finite trees is one
(to six decimal places).

%=========================================================================
\section{Estimating a distribution from counts: smoothing}\label{sec:smoothing}

Fix one context $u$. In the training corpus, nodes with context $u$ have symbols with counts
$c_{w}$ ($w$ one of the 13 symbols), $c = \sum_w c_w$ in total. How should we estimate $G_u$?

\paragraph{Maximum likelihood.} $\hat G_u(w) = c_w/c$. This gives probability zero to every symbol not
seen in context $u$, and so (by the product~\eqref{eq:product}) probability zero to every tree that
uses an unseen (context, symbol) pair. With hundreds of contexts and limited data this is fatal.

\paragraph{Dirichlet prior.} Put a Dirichlet prior on $G_u$ with mean $H$ (some distribution over the
symbols) and total concentration $\theta>0$, i.e.\ $G_u \sim \mathrm{Dirichlet}(\theta H)$. The
posterior predictive probability of the next symbol is
\begin{equation}\label{eq:dirichlet}
  P(w \mid \text{counts}) = \frac{c_w + \theta H(w)}{c + \theta}.
\end{equation}
This is a weighted average of the observed frequencies $c_w/c$ and the prior mean $H(w)$, with weights
$c/(c+\theta)$ and $\theta/(c+\theta)$. Every symbol now has positive probability. A small $\theta$ trusts
the counts; a large $\theta$ trusts $H$.

\paragraph{The Chinese restaurant picture.} Equation~\eqref{eq:dirichlet} has a well-known sequential
description. Picture a restaurant whose customers are the training observations, seated in order. Each
table serves one dish (a symbol). A new customer joins an existing table with probability proportional
to the number of people already at it, or opens a new table with probability proportional to $\theta$,
in which case the new table's dish is drawn from $H$. Summing over tables, the probability that the next
customer eats $w$ is exactly~\eqref{eq:dirichlet}. The table picture looks like extra baggage here, but
it is what makes the next two steps possible.

%=========================================================================
\section{The Pitman--Yor process}\label{sec:py}

The Pitman--Yor process $\PY(d, \theta, H)$ adds a \emph{discount} $0 \le d < 1$ to the restaurant:
\begin{itemize}
  \item join an existing table $k$ with probability proportional to (its size $-\ d$);
  \item open a new table with probability proportional to $\theta + d\,t$, where $t$ is the current
  number of tables, and draw its dish from $H$.
\end{itemize}
With $t_w$ the number of tables serving $w$ (and $t = \sum_w t_w$), the predictive probability of $w$ is
\begin{equation}\label{eq:py}
  P(w \mid \text{seating}) = \frac{c_w - d\,t_w + (\theta + d\,t)\, H(w)}{\theta + c}.
\end{equation}
With $d=0$ this is the Dirichlet case~\eqref{eq:dirichlet}. With $d>0$, a fixed amount $d$ is taken from
every table and given to the base $H$; the more distinct tables there are, the more mass goes to $H$.
This matches a pattern seen in language and in our corpus alike: many symbols (or forms) are rare, and
a discount moves mass from the frequently counted ones to the unseen ones more effectively than a
Dirichlet does. Both $d$ and $\theta$ are learned from the data (Section~\ref{sec:inference}).

%=========================================================================
\section{Hierarchical Pitman--Yor: backing off to shorter contexts}\label{sec:hpy}

What should the base $H$ be for context $u$? A natural choice is the distribution for a \emph{shorter}
context: if $u = (\texttt{*}, \texttt{a})$ (``second child of a product whose first child is a
constant''), its natural fallback is $(\texttt{a})$ or $(\texttt{*})$, which has more data. Write
$\pi(u)$ for the context obtained by dropping one component of $u$ (which one is dropped first is
fixed for each prior; by default the most distant ancestor goes first, as in \texttt{fnprior}). The
hierarchical Pitman--Yor model (Teh, 2006) is
\begin{align}
  G_u &\sim \PY\big(d_{g(u)},\ \theta_{g(u)},\ G_{\pi(u)}\big) \quad \text{for every context } u, \label{eq:hpy}\\
  G_{\emptyset} &\sim \PY\big(d_{0},\ \theta_{0},\ \Unif(13 \text{ symbols})\big), \notag
\end{align}
where $\emptyset$ is the empty context (all components dropped) and $g(u)$ is the \emph{group} of $u$:
contexts with the same set of components and the same child position (first or second) share one pair
$(d, \theta)$.

In the restaurant picture, there is one restaurant per context. Each training node is a customer in the
restaurant of its full context. Whenever a restaurant opens a new table, the dish for that table is drawn
from the parent distribution $G_{\pi(u)}$; in the restaurant picture, this sends one customer to the
parent restaurant. So the parent restaurant's customers are the \emph{tables} of its children, not the
raw observations.

The predictive probability is~\eqref{eq:py} applied recursively, with $H = $ the parent's predictive:
\begin{equation}\label{eq:hpy-pred}
  P(w \mid u) = \frac{c_{uw} - d\, t_{uw} + (\theta + d\, t_{u})\, P(w \mid \pi(u))}{\theta + c_u},
\end{equation}
ending in the uniform distribution. If no training node had context $u$ ($c_u=0$), this is simply the
parent's prediction. So a tree whose contexts were never seen still gets a positive probability: each
factor of~\eqref{eq:product} falls back to shorter and shorter contexts, and finally to the symbol
frequencies of the empty context.

\paragraph{Why table counts in the parent matter.} Suppose the symbol \texttt{x} appears 500 times, but
always as the second child of \texttt{+} after an \texttt{a}. The full-context restaurant handles those
500 well. The parent restaurant, which is used for \emph{other} contexts, receives only a few customers
(one per table), so it does not conclude that \texttt{x} is overwhelmingly likely everywhere. The
parent's counts measure in how many different situations a symbol is used, which is the right thing to
use for unfamiliar contexts.

\subsection{Inference}\label{sec:inference}

The seating arrangement (which customers share tables, and so the $t_{uw}$) is not observed, and neither
are the $d$'s and $\theta$'s. We sample them by MCMC:
\begin{itemize}
  \item \emph{Seating.} Gibbs sampling: remove each customer in turn and re-seat it from its conditional
  distribution given everyone else (join a table with probability $\propto$ size $-\ d$, open a new one
  with probability $\propto (\theta + d t)\times$ the parent's predictive; opening a table adds a
  customer to the parent, recursively).
  \item \emph{Hyperparameters.} Prior $d \sim \mathrm{Beta}(1,1)$ (uniform on $[0,1)$) and
  $\theta \sim \mathrm{Gamma}(1,1)$ for each group; random-walk Metropolis updates using the
  probability of the current seating.
\end{itemize}
We run 20 sweeps, discard the first 10, and keep a snapshot of the seating and hyperparameters after each
of the last 10. The prior probability of a tree is the average of~\eqref{eq:product} over the 10
snapshots, which approximates the posterior predictive probability given the corpus.

%=========================================================================
\section{The tree priors}\label{sec:variants}

All of these share the coding of Section~\ref{sec:trees} and so define distributions over exactly the
same set of trees; their probabilities can be compared directly.

\subsection{Baselines}

\paragraph{\texttt{uniform symbols}.} Every symbol at every node has probability $1/13$, so
$\log P(T) = -m \log 13$ for a tree with $m$ nodes. It learns nothing and is improper
(Section~\ref{sec:proper}); it is listed only as a reference point.

\paragraph{\texttt{sym0}.} The hierarchical Pitman--Yor model with $n=0$: the context contains no
symbols at all, only whether the node is a first child (or root) or a second child. So there are just two
distributions, both learned from the corpus: how often each symbol is used as a first argument and as a
second argument. A tree's probability then depends only on how many nodes it has and which symbols they
are, not on how the symbols are arranged. This is the honest baseline: any gain of a richer prior over
\texttt{sym0} comes from learning the structure of expressions, not just their vocabulary.

\subsection{Context models}

\paragraph{\texttt{hpy2}.} The hierarchical Pitman--Yor model with $n=2$, the contexts of
Example~\ref{ex:contexts}. First children (and the root) are conditioned on (grandparent, parent) and
back off to (parent) and then to the empty context; second children are conditioned on (parent, left
sibling) and back off to (left sibling) and then to the empty context. This is the fully Bayesian
counterpart of \texttt{fnprior} (next).

\paragraph{\texttt{hpy}.} An earlier, separate implementation of the same model as \texttt{hpy2}. The two
agree to within $0.01$ nats per form and are kept as a check.

\paragraph{\texttt{fnprior\_kn} and \texttt{fnprior\_kn\_tuned}.} The prior of Bartlett, Desmond \&
Ferreira (\texttt{fnprior} version~0.5.0), retrained on our corpus with their settings $n=2$,
a separate root context and constants coded \texttt{a}. It uses \emph{interpolated Kneser--Ney
smoothing}:
\[
  P(w \mid u) = \frac{\max(c_{uw} - D,\ 0)}{c_u} + \frac{D\, N_{1+}(u)}{c_u}\, P(w \mid \pi(u)),
\]
where $N_{1+}(u)$ is the number of distinct symbols seen in context $u$, $D=0.8$, and the lower-order
distributions are estimated from ``continuation counts'' (the number of distinct longer contexts a
symbol appears in) rather than raw counts. Comparing with~\eqref{eq:hpy-pred}: Kneser--Ney is exactly
the hierarchical Pitman--Yor predictive with $\theta = 0$, $d = D$ fixed, and exactly one table per
distinct (context, symbol) pair ($t_{uw} = 1$ whenever $c_{uw}>0$); the continuation counts are the
table counts. So \texttt{fnprior} is a fixed-parameter, single-seating approximation to
\texttt{hpy2}, and the two give nearly identical scores. \texttt{fnprior\_kn\_tuned} instead chooses $D$ by 10-fold
cross-validation on the training data (\texttt{fnprior}'s own tuning routine; one discount for first
children and one for second children). (\texttt{fnprior}'s
vocabulary also includes a second variable, \texttt{y}; the small mass it gives to \texttt{y} is mass
outside our one-variable space.)

\paragraph{\texttt{hpy3\_depth}.} As \texttt{hpy2}, but with longer and richer contexts: $n=3$
(a first child sees three ancestors, a second child two ancestors and its left sibling) and the node's
depth (capped at 6) as a further component. Backing off drops the depth first, then the most distant
ancestor, and so on. The depth lets the prior learn that expressions rarely nest deeply, so it matches
the size distribution of the corpus better than \texttt{hpy2}.

\paragraph{\texttt{hpy2\_src}.} As \texttt{hpy2}, with the corpus source $s$ of each training form as an
extra, outermost context component: each source has its own restaurants, which back off to restaurants
shared by all sources. This lets sources differ in style while sharing strength. A form from a source
not in the training data is scored with the shared level, which is the correct prediction for a new
source (a new source's own restaurant would be empty, and an empty restaurant predicts with its parent).

\subsection{Adaptor priors: reusing whole forms and subtrees}\label{sec:adaptors}

The context models build every tree from scratch, one symbol at a time. But the corpus contains the same
forms over and over: Michaelis--Menten, the Hill function, the exponential. A context model can only
give Michaelis--Menten the product of its seven symbol probabilities, which is much smaller than its
actual frequency in the corpus. An \emph{adaptor} (Johnson, Griffiths \& Goldwater, 2007) adds a second
layer that can reuse whole trees.

At an adapted position $u$ the tree is drawn from
\[
  H_u \sim \PY\big(a,\ b,\ B_u\big), \qquad
  B_u(T) = G_u(\text{root symbol of } T) \times \prod_{\text{children } i} P_{u_i}(T_i),
\]
where $G_u$ is a context model of the previous section (the \emph{base}), $u_i$ is the position of
child $i$, and $P_{u_i}$ is that position's own adaptor $H_{u_i}$ if it is adapted, or the base
$B_{u_i}$ if not. In the restaurant picture, the customers are now whole (sub)trees and each table
serves a whole tree. A new tree either joins a table serving the same tree (\emph{reusing} it whole), or
opens a new table, in which case it is built symbol by symbol from the base. The predictive probability
of a tree $T$ at an adapted position is
\begin{equation}\label{eq:adaptor}
  P(T) = \frac{n_T - a\, K_T + (b + a K)\, B_u(T)}{b + n},
\end{equation}
where $n$ is the number of trees seen at that position, $K$ the number of tables, and $n_T$, $K_T$ the
customers and tables of tree $T$.

The two terms of~\eqref{eq:adaptor} explain the behaviour we see in the results. A form seen often gets
roughly its corpus frequency $(n_T - aK_T)/(b+n)$, far more than the base alone would give. A form never
seen gets only the base probability multiplied by the share kept for new tables,
$(b + aK)/(b + n)$, which is well below one when the corpus repeats a modest number of forms many
times. So adaptors win on familiar forms and pay a fixed toll on new ones.

The table counts also change what the base learns: the base sees one customer per \emph{table}, roughly
one per distinct form, rather than one per occurrence. So frequent forms do not dominate the base, which
then models the variety of forms rather than their popularity.

The adaptor's parameters $(a, b)$ are shared within three groups (the root, first children, second
children), with priors $a \sim \mathrm{Beta}(1,1)$ and $b \sim \mathrm{Gamma}(1,1)$. Inference is again
MCMC: each training tree's seating is removed and a new one proposed from the top down, accepted or
rejected by a Metropolis--Hastings step (exact, because the proposal's probability is tracked); between
these updates the base's seating and all hyperparameters are updated as in
Section~\ref{sec:inference}. Predictions average over 10 snapshots, as before.

\paragraph{\texttt{form\_ad}.} Only the root is adapted, with an \texttt{hpy2} base: a Pitman--Yor
distribution over whole forms that falls back on \texttt{hpy2} for new forms.

\paragraph{\texttt{form\_ad3\_depth}.} As \texttt{form\_ad}, with the \texttt{hpy3\_depth} base.

\paragraph{\texttt{tree\_ad3}.} Every position is adapted (any subtree can be reused anywhere, e.g.\ a
saturating term $x/(c+x)$ inside a larger form), with an $n=3$ base without depth.

%=========================================================================
\section{Grammar priors (PCFGs)}\label{sec:pcfg}

The second family generates forms with a probabilistic context-free grammar (PCFG) instead of the
\texttt{fnprior} tree coding. The grammar is
\begin{align*}
  E &\to E + F \mid E - F \mid F && \text{(sums)}\\
  F &\to F \times T \mid F / T \mid T && \text{(products)}\\
  T &\to (E) \mid R(E) \mid V \mid C \mid (E)^{C} && \text{(factors)}\\
  R &\to \exp \mid \log \mid \sin \mid \mathrm{sqrt} \mid \tanh && \text{(functions)}
\end{align*}
with $V$ the variable and $C$ a constant. A form is generated by starting from $E$ and repeatedly
replacing the left-most nonterminal by one of its options; its probability is the product of the
probabilities of the options chosen. Each form is scored by one canonical derivation.

For each nonterminal $N$ and context $u$ the option probabilities have a Dirichlet prior,
$p_{u,N} \sim \mathrm{Dirichlet}(\alpha_N)$, with fixed $\alpha$ (e.g.\ $\alpha_E = (1,1,4)$, favouring
$E \to F$). Given the counts of each option in the training derivations, the posterior mean probability
is~\eqref{eq:dirichlet}. The schemes differ only in the context $u$:
\begin{itemize}
  \item \texttt{default\_prior}: no training at all; the prior mean of the Dirichlet, no context;
  \item \texttt{cf}: context-free (one distribution per nonterminal), trained;
  \item \texttt{parent}: the context is the nearest enclosing function (none, $\exp$, $\log$, \dots, or
  the base of a power);
  \item \texttt{position}: as \texttt{parent}, plus whether a factor is in a numerator or a denominator;
  \item \texttt{term}: as \texttt{position}, plus whether a product is one of several terms of a sum;
  \item \texttt{term\_root}: as \texttt{term}, plus whether the choice is made at the top level of the
  expression (a start-of-tree context, like \texttt{fnprior}'s root context).
\end{itemize}
A context unseen in training uses the prior mean. In the \texttt{\_backoff20} versions each context's
Dirichlet is centred on the trained context-free fit instead,
$\alpha_{u,N} = 20\, \hat p_{\texttt{cf},N} + 0.1$, a one-level version of the back-off idea of
Section~\ref{sec:hpy}.

Unlike the tree priors, the grammar has no discount, no learned hyperparameters and only one level of
back-off, so it cannot share strength as flexibly. It also generates a different set of objects
(derivations, in which some different strings denote the same tree), so its raw probabilities are only
approximately comparable with the tree priors'.

%=========================================================================
\section{Training data and how the priors are compared}\label{sec:evaluation}

\paragraph{Training data.} Each source is a list of (model, form) pairs, each model contributing each
of its distinct forms once. Sources differ greatly in size, so we train in two ways: \emph{pooled}
(every pair counts once, so large sources dominate) and \emph{equal} (each source is replicated to the
size of the largest, so each source carries the same total weight).

\paragraph{Score.} A prior is scored on test forms it was not trained on by its own log probability
$\log P(\text{form})$, averaged over test (model, form) pairs; higher is better and the unit is nats.
Because the tree priors are proper distributions over the same space, a difference of $1$ nat means the
better prior gives the test forms, on average, $e \approx 2.7$ times more probability. We do not
renormalise over a finite list of candidate forms: renormalising rewards priors that are flat over that
list, which is an artefact of the list.

\paragraph{Seen and novel forms.} A test form is \emph{seen} if its key occurs in the training data and
\emph{novel} otherwise. The two kinds test different things: seen forms test how well a prior
reproduces the corpus's habits; novel forms test how sensibly it extrapolates. Novel forms are also
larger on average, and every proper prior must give larger forms less probability, so comparisons of
novel forms are made relative to \texttt{sym0} or at fixed tree size.

\paragraph{Test designs.} Leave one source out (train on all sources but one, test on that one),
hold-out sources (train on all twelve training sources, test on three sources never used for training),
and, in preparation, a temporal split (train on models published before a cutoff year, test on those
published after).

%=========================================================================
\section*{References}
\begin{itemize}
  \item D.~J.~Bartlett, H.~Desmond and P.~G.~Ferreira (2023). Priors for symbolic regression.
  arXiv:2304.06333.
  \item M.~Johnson, T.~L.~Griffiths and S.~Goldwater (2007). Adaptor grammars: a framework for
  specifying compositional nonparametric Bayesian models. \emph{Advances in Neural Information
  Processing Systems} 19.
  \item S.~Goldwater, T.~L.~Griffiths and M.~Johnson (2006). Interpolating between types and tokens by
  estimating power-law generators. \emph{Advances in Neural Information Processing Systems} 18.
  \item Y.~W.~Teh (2006). A hierarchical Bayesian language model based on Pitman--Yor processes.
  \emph{Proceedings of COLING/ACL}, 985--992.
\end{itemize}

\end{document}
